% ============================================================================
%  ВВЕДЕНИЕ
% ============================================================================

\introduction

% ---------------------------------------------------------------------------
\actuality

Промышленные музыкальные сервисы сегодня оперируют каталогами, в которые ежедневно добавляются десятки тысяч новых треков. Только Spotify к 2024 году преодолел отметку в 100 миллионов композиций. Темпы роста каталогов давно превысили физические возможности ручной разметки. В таких условиях возрастает спрос на автоматические инструменты семантической разметки аудио: определение жанра, оценка эмоционального содержания, выделение структурных секций. На рынке такие задачи в большинстве случаев решаются специализированными моделями, каждая под свою задачу, что приводит к избыточному дублированию параметров и не использует общую акустическую природу входных данных. Настоящая работа направлена на проверку гипотезы о том, что объединение четырёх задач (структурной сегментации, классификации возбуждения, валентности и жанра) в единой мультизадачной модели с общим энкодером и независимыми головами не приводит к существенной деградации качества по сравнению с однозадачными моделями, и при этом даёт согласованную разметку за один проход.

% ---------------------------------------------------------------------------
\review

Предметная область опирается на обзорные работы по MIR~\cite{downie2003mir} и применению глубокого обучения~\cite{choi2017tutorial}, на исследования по структурной сегментации (метод матриц автоподобия, глубокие модели, датасет Harmonix~Set), распознаванию эмоций (модель Расселла, датасет DEAM) и жанровой классификации (бенчмарк GTZAN, набор моделей PANNs~\cite{kong2020panns} и архитектура AST~\cite{gong2021ast}). Эти источники служат для формирования гипотез и обоснования сравниваемых архитектур. Полный обзор приведён в Главе~1.

% ---------------------------------------------------------------------------
\goalandtasks

% ---------------------------------------------------------------------------
\tools

Работа опирается на стандартные для области методы цифровой обработки сигналов (короткое преобразование Фурье, лог-мел-спектрограмма, chroma), машинного обучения (свёрточные сети с SE-блоками, двунаправленные LSTM, самовнимание AST с трансфером весов AudioSet), регуляризации (SpecAugment, focal loss, label smoothing, MixUp) и динамического программирования (декодер Витерби с позиционными приорами). Реализация выполнена на Python 3.12 с PyTorch, librosa и scikit-learn. Веб-прототип построен на React 19 с Vite на фронтенде, на Elysia с Bun на бэкенде и на FastAPI с PyTorch на стороне инференса. Полноэкранный визуализатор Vibe Mode использует WebGL2.

% ---------------------------------------------------------------------------
\results

% ---------------------------------------------------------------------------
\originality

Научная новизна работы заключается в следующем. Разработана мультизадачная нейросетевая модель с единым аудио-энкодером и четырьмя классификационными головами, обучаемая одновременно на трёх разнородных датасетах (DEAM, Harmonix Set, GTZAN) по схеме round-robin. Впервые, насколько известно автору, выполнено систематическое сравнение применимости пяти конфигураций (CNN per-window, CNN с надстройкой BiLSTM, PANNs с линейными головами, PANNs с BiLSTM, Audio Spectrogram Transformer с частичной разморозкой) к этой задаче в едином экспериментальном протоколе. Предложен воспроизводимый пайплайн кросс-датасетного обучения через RoundRobinLoader с покомпонентным взвешиванием потерь, в том числе для случая, когда метки каждой задачи присутствуют только в одном из источников. Дополнительно предложен и количественно оценен составной постобработчик сегментных предсказаний (декодер Витерби с матрицей переходов, позиционными приорами и поклассовыми ограничениями минимальной длительности), повышающий Boundary~F1 на 14 пунктов относительно базового аргмакса и сокращающий среднее число выдаваемых сегментов с 28 до 7-8. Также реализован интерактивный веб-прототип с шестипанельным синхронизированным дашбордом и WebGL2-визуализатором, обеспечивающий академически честную коммуникацию ограничений модели (топ-3 жанров с предупреждением об уверенности при out-of-distribution).

% ---------------------------------------------------------------------------
\approbation

Результаты исследования прошли апробацию на студенческой научной конференции факультета прикладной математики и кибернетики Тверского государственного университета (2026 год). По теме работы принята к публикации статья «Мультизадачный анализ музыкальных композиций с помощью предобученного аудио-трансформера» в сборнике трудов конференции. В статью включены лучшие по каждой задаче результаты среди исследованных конфигураций. Помимо этого, разработанный прототип был отдельно протестирован на наборе из 41 аудиотрека различных жанров. Для каждого модель сформировала структурную разметку, эмоциональные кривые и распределение жанров. Реализация оформлена в виде открытого репозитория с 128 автоматическими тестами, обеспечивающими воспроизводимость экспериментов.

% ---------------------------------------------------------------------------
\structure

Работа состоит из введения, семи глав основной части, заключения, списка использованных источников и пяти приложений.

Глава~1 содержит обзор предметной области. В ней разбираются три основные задачи MIR из числа решаемых в работе (структурная сегментация, распознавание эмоций, жанровая классификация), эталонные датасеты по каждой задаче и пять архитектурных семейств для аудио-классификации: свёрточные сети с блоками канального внимания, надстройки двунаправленной LSTM, замороженные эмбеддинги PANNs CNN14 и Audio Spectrogram Transformer с трансфером с AudioSet. Отдельные параграфы посвящены мультизадачному обучению и техникам аккуратного дообучения предобученных моделей.

Глава~2 содержит формальную постановку задачи. Цель работы и четыре класса выходных предсказаний модели формализованы математически: входной аудиосигнал, последовательность лог-мел окон, четыре распределения вероятностей и постобработка во временные интервалы. На основе этой формализации сформулированы пять исследовательских вопросов RQ1--RQ5, легших в основу экспериментальной программы. Завершает главу содержательное обоснование выбора мультизадачного подхода по трём линиям: семантической, методологической и практической.

Глава~3 описывает данные. Подробно разбираются три открытых датасета (DEAM, Harmonix Set, GTZAN), включая их объёмы, лицензии и процедуру скачивания недостающего аудио для Harmonix через \texttt{yt-dlp}. Описан унифицированный пайплайн извлечения логарифмических мел-спектрограмм, окновая нарезка с шагом одна секунда, глобальная нормализация по обучающему сплиту и стратегия разбиения на обучающую, валидационную и тестовую выборки на уровне треков. Отдельный раздел посвящён компенсации сильного дисбаланса классов в структурной задаче через сочетание focal loss, весов классов с верхним порогом и поканальных весов потери.

Глава~4 содержит описание методов. Подробно разобраны все пять архитектурных конфигураций (CNN, CNN+BiLSTM, PANNs+Linear, PANNs+BiLSTM, AST~v2), общая схема мультизадачных голов и функция потерь на основе focal loss, оптимизатор AdamW с расписанием CosineAnnealingWarmRestarts. Отдельный большой раздел разбирает постобработку сегментных предсказаний на основе декодера Витерби: матрица переходов, позиционные приоры, поклассовые ограничения минимальной длительности, расширенная версия v2 с шестью механизмами. Завершает главу описание пайплайна инференса на пользовательских записях.

Глава~5 представляет результаты экспериментов. Описаны протокол оценки и гиперпараметры обучения, после чего излагаются четыре последовательных эксперимента: базовый запуск v1, улучшенный v2, расширение трансформер-архитектурой AST~v2 и неудачный CNN~v3. Сводное сравнение всех конфигураций даётся в отдельном разделе вместе с ablation-экспериментом по компонентам постобработки.

Глава~6 описывает реализованный веб-прототип. Разобрана трёхзвенная архитектура (фронтенд, бэкенд, сервис инференса), синхронизированный дашборд из шести панелей, ключевые UX-решения для честной коммуникации ограничений модели и технологический стек. Завершает главу описание визуального редизайна в стилистике glassmorphism и полноэкранного визуализатора Vibe Mode на WebGL2-шейдере.

Глава~7 содержит сводный анализ результатов. Даны ответы на все пять исследовательских вопросов, зафиксирована финальная конфигурация прототипа, систематически проанализированы пять групп ограничений подхода с возможными направлениями развития, проведено сопоставление со state-of-the-art по каждой задаче, сформулирована научная новизна и практическая значимость работы.

Заключение кратко перечисляет основные результаты и направления дальнейшей работы. Список использованных источников содержит 40 наименований. Пять приложений (А--Д) собирают подробные технические сведения: характеристики датасетов, гиперпараметры конфигураций, путеводитель по графическим результатам, JSON-схему выхода модели и команды воспроизведения экспериментов.
